%% ══════════════════════════════════════════════════════════════
%%  CHAPTER 7 — MODELING & SCIENTIFIC LAYER
%% ══════════════════════════════════════════════════════════════
\chapter{Modeling and Scientific Computation Layer}
\label{ch:modeling}

\section{Overview}
\label{sec:modeling:overview}

The scientific computation layer is the analytical engine of \autoscholar, responsible for all quantitative analysis, model estimation, and figure generation.
It consists of a sandboxed Python execution environment with 17 pre-installed scientific libraries, 7 dedicated analysis services, and a custom figure generation pipeline.
All computation runs in short-lived Python subprocesses spawned by TypeScript bridge modules.

\section{Python Execution Architecture}
\label{sec:modeling:python}

\subsection{The TypeScript--Python Bridge Pattern}

Every Python computation follows a consistent architectural pattern:

\begin{enumerate}
  \item A TypeScript module (\file{*.ts}) defines input/output types, validates parameters, and generates displayable code.
  \item The bridge spawns a Python subprocess (\file{*.py}) via \texttt{child\_process.spawn}.
  \item Input is passed as JSON via stdin.
  \item Python executes the analysis and returns structured JSON via stdout.
  \item The TypeScript module parses the output and returns typed results.
\end{enumerate}

This pattern is replicated across 10+ service pairs:

\begin{table}[H]
\centering
\caption{TypeScript--Python service pairs.}
\label{tab:ts-py-pairs}
\begin{tabular}{lll}
\toprule
\textbf{TypeScript Bridge} & \textbf{Python Service} & \textbf{Function} \\
\midrule
\file{garch.ts} & \file{garchService.py} & GARCH volatility modeling \\
\file{varCalculator.ts} & \file{varService.py} & Value at Risk \\
\file{monteCarlo.ts} & \file{monteCarloService.py} & Monte Carlo simulation \\
\file{portfolioOptimizer.ts} & \file{portfolioOptimizer.py} & Portfolio optimization \\
\file{riskMetrics.ts} & \file{riskMetrics.py} & Risk/performance metrics \\
\file{optionsPricing.ts} & \file{optionsPricingService.py} & Black-Scholes pricing \\
\file{volatilitySurface.ts} & \file{volatilitySurfaceService.py} & IV surface construction \\
\file{plotExecutor.ts} & \file{plotService.py} & Chart generation \\
\file{dataDownload.ts} & \file{dataDownloadService.py} & Data export \\
\file{customPython.ts} & \file{customPythonExecutor.py} & Arbitrary analysis \\
\bottomrule
\end{tabular}
\end{table}

\subsection{Custom Python Executor}

The custom executor (\file{customPythonExecutor.py}, 1,774 lines) is the most powerful and complex Python service.
It provides an unrestricted execution environment for Claude-generated Python code with:

\subsubsection{Pre-Installed Libraries}

\begin{multicols}{3}
\begin{itemize}
  \item NumPy
  \item Pandas
  \item SciPy
  \item Matplotlib
  \item Scikit-learn
  \item Statsmodels
  \item XGBoost
  \item LightGBM
  \item ARCH (for GARCH)
  \item CVXPY
  \item Seaborn
  \item Plotly
  \item OpenPyXL
  \item python-docx
  \item DuckDB
  \item Requests
  \item BeautifulSoup
\end{itemize}
\end{multicols}

\subsubsection{Utility Functions}

\begin{description}
  \item[\func{save\_excel}] Creates styled multi-sheet XLSX workbooks with title rows, auto-column widths, conditional formatting (color scales, data bars), summary statistics (sum, average, min, max, count), and formula support.
  \item[\func{save\_word}] Generates formatted DOCX documents with styled headings, tables, and paragraphs.
  \item[\func{df\_to\_latex}] Converts DataFrames to \LaTeX{} table markup with proper formatting.
  \item[\func{save\_latex}] Compiles \LaTeX{} content into PDF with embedded figures.
\end{description}

\subsubsection{Automatic Error Patching}

The TypeScript bridge (\file{customPython.ts}, 371 lines) applies automatic fixes before execution:

\begin{table}[H]
\centering
\caption{Automatic Python code patches.}
\label{tab:python-patches}
\begin{tabularx}{\textwidth}{lX}
\toprule
\textbf{Pattern} & \textbf{Fix} \\
\midrule
Deprecated pandas frequencies & \texttt{'M'} $\to$ \texttt{'ME'}, \texttt{'Q'} $\to$ \texttt{'QE'}, \texttt{'Y'} $\to$ \texttt{'YE'} \\
NumPy deprecations & \texttt{np.trapz} $\to$ \texttt{np.trapezoid}, remove \texttt{np.math} \\
Scikit-learn changes & t-SNE \texttt{n\_iter} $\to$ \texttt{max\_iter} \\
Matplotlib conflicts & Monkey-patch \texttt{tight\_layout} for colorbar compatibility \\
fillna deprecation & \texttt{fillna(method='ffill')} $\to$ \texttt{ffill()} \\
Data fabrication & Flag random data generation without API calls \\
\bottomrule
\end{tabularx}
\end{table}

\section{Dedicated Analysis Services}
\label{sec:modeling:services}

\subsection{GARCH Volatility Modeling}

\file{garchService.py} (224 lines) implements GARCH($p$,$q$) volatility estimation:

\begin{equation}
\sigma_t^2 = \omega + \sum_{i=1}^{p} \alpha_i \varepsilon_{t-i}^2 + \sum_{j=1}^{q} \beta_j \sigma_{t-j}^2
\label{eq:garch}
\end{equation}

\subsubsection{Implementation Details}

\begin{enumerate}
  \item Fetches historical prices from FMP API.
  \item Calculates percentage returns: $r_t = 100 \times \frac{P_t - P_{t-1}}{P_{t-1}}$.
  \item Fits a constant-mean GARCH model with normal distribution using the \pkg{arch} library.
  \item Produces:
  \begin{itemize}
    \item Conditional volatility time series.
    \item Volatility forecasts (1--252 days ahead).
    \item Model parameters ($\omega$, $\alpha_i$, $\beta_j$).
    \item Persistence: $\sum \alpha_i + \sum \beta_j$.
    \item AIC/BIC for model comparison.
    \item Interpretive labels (mean-reverting, shock impact, clustering strength).
  \end{itemize}
\end{enumerate}

\subsubsection{Validation}

Input validation (\file{garch.ts}, 282 lines):
\begin{itemize}
  \item $p, q \in [1, 5]$.
  \item Forecast horizon $\in [1, 252]$ days.
  \item Data period $\in [30, 2520]$ days ($\approx$ 10 years max).
\end{itemize}

\subsection{Value at Risk (VaR)}

\file{varService.py} (307 lines) implements four VaR methodologies:

\subsubsection{Historical VaR}

\begin{equation}
\text{VaR}_\alpha^{\text{hist}} = -\text{Quantile}(r_1, r_2, \ldots, r_T; \alpha)
\end{equation}

Uses the empirical return distribution directly.

\subsubsection{Parametric VaR}

\begin{equation}
\text{VaR}_\alpha^{\text{param}} = -(\mu + z_\alpha \sigma)
\end{equation}

where $z_\alpha$ is the standard normal quantile.
Assumes normally distributed returns.

\subsubsection{Monte Carlo VaR}

\begin{equation}
\text{VaR}_\alpha^{\text{MC}} = -\text{Quantile}\left(\{r_i^{\text{sim}}\}_{i=1}^{N}; \alpha\right)
\end{equation}

where $r_i^{\text{sim}} \sim \mathcal{N}(\hat{\mu}, \hat{\sigma}^2)$ with $N$ simulations (default: 10,000).

\subsubsection{Conditional VaR (Expected Shortfall)}

\begin{equation}
\text{CVaR}_\alpha = \mathbb{E}\left[r \mid r < -\text{VaR}_\alpha\right]
\end{equation}

The average loss beyond the VaR threshold.

\subsubsection{Outputs}

Results are computed at three confidence levels (90\%, 95\%, 99\%) with both percentage and dollar amounts.
Distribution statistics include skewness, kurtosis, and Jarque-Bera normality test.
The system recommends the appropriate VaR method based on distribution characteristics.

\subsection{Monte Carlo Simulation}

\file{monteCarloService.py} (152 lines) performs investment projections:

\begin{equation}
P_{t+1} = P_t \times (1 + r_t), \quad r_t \sim \mathcal{N}(\hat{\mu}_{\text{daily}}, \hat{\sigma}_{\text{daily}}^2)
\end{equation}

\subsubsection{Features}

\begin{itemize}
  \item Two modes: symbol-based (auto-fetches from FMP) and manual (user-provided prices).
  \item Configurable simulations (1--100,000), time horizon, and initial investment.
  \item Output: mean/median/percentile terminal values, probability of profit, sample paths, distribution data.
\end{itemize}

\subsection{Portfolio Optimization}

\file{portfolioOptimizer.py} (434 lines) implements Markowitz Mean-Variance Optimization:

\begin{equation}
\max_{\mathbf{w}} \frac{\mathbf{w}^T \boldsymbol{\mu} - r_f}{\sqrt{\mathbf{w}^T \boldsymbol{\Sigma} \mathbf{w}}}
\quad \text{s.t.} \quad \sum_i w_i = 1, \quad w_i \in [l_i, u_i]
\label{eq:mvo}
\end{equation}

\subsubsection{Outputs}

\begin{itemize}
  \item Maximum Sharpe ratio portfolio weights.
  \item Minimum volatility portfolio weights.
  \item 50-point efficient frontier.
  \item Correlation matrix.
  \item Comparison against equal-weight benchmark.
\end{itemize}

Uses SciPy's SLSQP optimizer for constrained optimization.

\subsection{Risk Metrics}

\file{riskMetrics.py} (543 lines) computes 20+ performance and risk metrics:

\begin{longtable}{lp{8cm}}
\caption{Risk and performance metrics.} \label{tab:risk-metrics} \\
\toprule
\textbf{Metric} & \textbf{Formula / Description} \\
\midrule
\endfirsthead
\toprule
\textbf{Metric} & \textbf{Formula / Description} \\
\midrule
\endhead
Annualized Return & $(1 + \bar{r}_{\text{daily}})^{252} - 1$ \\
Annualized Volatility & $\sigma_{\text{daily}} \times \sqrt{252}$ \\
Sharpe Ratio & $\frac{R_p - R_f}{\sigma_p}$ \\
Sortino Ratio & $\frac{R_p - R_f}{\sigma_{\text{downside}}}$ \\
Calmar Ratio & $\frac{R_p}{\text{MaxDrawdown}}$ \\
Information Ratio & $\frac{R_p - R_b}{\sigma(R_p - R_b)}$ \\
Treynor Ratio & $\frac{R_p - R_f}{\beta_p}$ \\
Beta & $\frac{\text{Cov}(R_p, R_m)}{\text{Var}(R_m)}$ \\
Alpha (Jensen's) & $R_p - [R_f + \beta(R_m - R_f)]$ \\
$R^2$ & From OLS of $R_p$ on $R_m$ \\
Max Drawdown & $\max_t \frac{\text{Peak}_t - \text{Value}_t}{\text{Peak}_t}$ \\
Recovery Days & Days from max drawdown to recovery \\
Tracking Error & $\sigma(R_p - R_b)$ \\
Upside Capture & $\frac{\bar{R}_p | R_m > 0}{\bar{R}_m | R_m > 0}$ \\
Downside Capture & $\frac{\bar{R}_p | R_m < 0}{\bar{R}_m | R_m < 0}$ \\
Win Rate & Fraction of positive return days \\
Profit Factor & $\frac{\sum \text{gains}}{\sum |\text{losses}|}$ \\
\bottomrule
\end{longtable}

Uses SPY as the default benchmark with date-aligned return series.

\subsection{Options Pricing}

\file{optionsPricingService.py} (238 lines) implements Black-Scholes for European options:

\begin{align}
C &= S \Phi(d_1) - K e^{-rT} \Phi(d_2) \label{eq:bs-call} \\
P &= K e^{-rT} \Phi(-d_2) - S \Phi(-d_1) \label{eq:bs-put} \\
d_1 &= \frac{\ln(S/K) + (r + \sigma^2/2)T}{\sigma\sqrt{T}} \\
d_2 &= d_1 - \sigma\sqrt{T}
\end{align}

\subsubsection{Greeks}

\begin{align}
\Delta_C &= \Phi(d_1), \quad \Delta_P = \Phi(d_1) - 1 \\
\Gamma &= \frac{\phi(d_1)}{S\sigma\sqrt{T}} \\
\Theta_C &= -\frac{S\phi(d_1)\sigma}{2\sqrt{T}} - rKe^{-rT}\Phi(d_2) \\
\mathcal{V} &= S\sqrt{T}\phi(d_1) \\
\rho_C &= KTe^{-rT}\Phi(d_2)
\end{align}

\subsubsection{Implied Volatility}

Computed via Newton-Raphson iteration:

\begin{equation}
\sigma_{n+1} = \sigma_n - \frac{C(\sigma_n) - C_{\text{market}}}{\mathcal{V}(\sigma_n)}
\end{equation}

\subsection{Volatility Surface}

\file{volatilitySurfaceService.py} (378 lines) constructs 3D implied volatility surfaces:

\begin{enumerate}
  \item Fetches option chain from Financial Data API.
  \item Intelligently samples up to 60 contracts across strikes and expirations.
  \item Retrieves implied volatility for each contract.
  \item Interpolates a smooth surface using SciPy's cubic/linear \texttt{griddata}.
  \item Renders a 3D matplotlib plot with scatter points for actual data.
\end{enumerate}

\section{Scientific Computing Libraries}
\label{sec:modeling:libraries}

\subsection{Econometric Modeling with Statsmodels}

Used for:
\begin{itemize}
  \item OLS, WLS, GLS regression.
  \item Panel data models (fixed effects, random effects).
  \item Time series models (ARIMA, VAR, VECM).
  \item Hypothesis testing (t-tests, F-tests, Jarque-Bera, ADF, KPSS).
  \item Robust standard errors (Newey-West, clustered).
\end{itemize}

\subsection{Machine Learning with Scikit-learn}

Used for:
\begin{itemize}
  \item Classification and regression models.
  \item Cross-validation and hyperparameter search.
  \item Feature selection and dimensionality reduction.
  \item Clustering and anomaly detection.
  \item SHAP value integration for interpretability.
\end{itemize}

\subsection{Gradient Boosting with XGBoost/LightGBM}

Used for:
\begin{itemize}
  \item High-performance predictive modeling.
  \item Feature importance ranking.
  \item Early stopping and regularization.
\end{itemize}

\subsection{Convex Optimization with CVXPY}

Used for:
\begin{itemize}
  \item Portfolio optimization beyond mean-variance (CVaR, risk parity).
  \item Robust optimization problems.
  \item Constrained optimization with complex constraints.
\end{itemize}

\subsection{Volatility Modeling with ARCH}

Used for:
\begin{itemize}
  \item GARCH, EGARCH, GJR-GARCH models.
  \item Multivariate GARCH (DCC, BEKK).
  \item Volatility forecasting.
\end{itemize}

\section{Figure Generation Pipeline}
\label{sec:modeling:figures}

\subsection{Style Configuration}

\file{matplotlib\_config.py} (79 lines) sets the publication-quality defaults:

\begin{itemize}
  \item \textbf{Fonts:} DejaVu Serif, Times New Roman (serif family).
  \item \textbf{Colors:} 8-color muted palette (navy, crimson, forest green, brown, purple, teal, olive, charcoal).
  \item \textbf{Layout:} Inward ticks, no top/right spines, print-safe grid (alpha 0.25).
  \item \textbf{Resolution:} 150 DPI for display, 300 DPI for save.
  \item \textbf{Layout engine:} Constrained layout (avoids \texttt{tight\_layout} conflicts).
\end{itemize}

\subsection{Figure Capture}

\file{figure\_capture.py} (104 lines) processes all open matplotlib figures:

\begin{enumerate}
  \item Iterates over \texttt{plt.get\_fignums()}.
  \item For each figure, checks if it is empty using two-stage analysis:
  \begin{itemize}
    \item \textbf{Stage 1:} Inspects axes for artists (lines, patches, images, collections, texts, tables).
    \item \textbf{Stage 2:} Renders to image and measures pixel variance (threshold: 50.0).
  \end{itemize}
  \item Non-empty figures receive the AutoScholar watermark (teal, bottom-right).
  \item Figures are encoded as base64 PNG strings.
  \item All figures are closed to free memory.
\end{enumerate}

\subsection{Vision Analysis}

\file{figureAnalysisService.ts} (113 lines) optionally sends figures to Claude Opus for vision-based analysis:

\begin{itemize}
  \item Chart type identification.
  \item Data description.
  \item Key trend identification.
  \item Notable value extraction.
  \item Parallel processing via \texttt{Promise.all}.
\end{itemize}

\section{Script Management and Repair}
\label{sec:modeling:repair}

\subsection{Script Lifecycle}

The ScriptManager (\file{scriptManager.ts}, 276 lines) tracks Python scripts through their lifecycle:

\begin{enumerate}
  \item \textbf{Save:} Script content saved to session-scoped temp directory with unique ID.
  \item \textbf{Execute:} Script run via the Python subprocess.
  \item \textbf{Error capture:} Structured error parsing via the StructuredLogger.
  \item \textbf{Repair:} Claude reads the script and errors, generates a patch.
  \item \textbf{Patch:} DiffEngine applies the patch (up to 8 attempts).
  \item \textbf{Re-run:} Patched script re-executed.
  \item \textbf{Cleanup:} Session scripts cleaned up on session end or after 24 hours.
\end{enumerate}

\subsection{Diff Engine}

The DiffEngine (\file{diffEngine.ts}, 400 lines) supports three patching strategies:

\begin{description}
  \item[Strategy 1: Standard unified diff] with fuzzy matching (up to 50-line offset tolerance). Parses \texttt{@@} hunk headers and applies line-by-line.
  \item[Strategy 2: Search-replace blocks] in three LLM-friendly formats:
  \begin{itemize}
    \item Marker-based: \texttt{<<<< SEARCH / ==== / >>>> REPLACE}
    \item Labeled format with code blocks.
    \item Raw \texttt{-/+} line extraction.
  \end{itemize}
  \item[Strategy 3: Trimmed/stripped matching] as a last-resort fallback that ignores whitespace differences.
\end{description}

\subsection{Structured Error Parsing}

The StructuredLogger (\file{structuredLogger.ts}, 145 lines) converts Python tracebacks into structured JSON:

\begin{lstlisting}[style=typescript, caption=Structured error output format.]
interface StructuredError {
  type: string;      // "KeyError", "ImportError", etc.
  message: string;
  file: string;
  line: number;
  traceback: string;
}
\end{lstlisting}

This structured format enables the \governor to classify errors by type and generate pattern-specific recovery advice.
